\section*{Introduction}
Our brains are primed to process faces, reflecting the central role faces play in our social lives, including conveying important information about identity, intent, and affect \citep{johnson_newborns_1991, palmer_face_2020, powell_social_2018}.
Emotional expressions are particularly consequential as they enable observers to infer others' mental states and modify behaviour accordingly in real time, which are essential for effective social interactions. 
Until recently, nearly all social interactions were between two (or more) people, however, that is changing. 
In recent years, there has been a proliferation of the use of avatars, (computer-generated human representations) across social media, gaming, virtual reality, and augmented reality.
As interaction with avatars eventually becomes commonplace, the capacity for us to accurately process the emotional expressions on virtual faces is essential for evaluating their utility and social acceptance in contexts requiring interactions \citep{kegel_dynamic_2020}. 
Despite a large body of work examining the neural mechanisms associated with processing emotions on real faces \citep{barrett_are_2006}, the neural systems supporting the recognition and interpretation of emotions on virtual faces is unclear. 
For instance, are emotional expressions in virtual faces processed differently from processing emotional expressions in real faces? 
Do emotional expressions on virtual faces engage the same neural mechanisms as they do on real faces? 

Emotion recognition on real faces has been driven by two models: the first of which is Basic Emotion Theory (BET), which posits the existence of a set of basic emotions \citep{ekman1971constants}. 
These emotions consist of Joy, Sadness, Anger, Fear, Disgust, and Surprise, with Neutral serving as a baseline. 
These basic emotions are reliably associated with distinct facial configurations, known as the Facial Action Coding System (FACS), which catalogues the facial muscle actions that produce visible expressions \citep{ekman1978facial}.
Cross-cultural research has long supported the biological roots and broad universality of these expressions, even preliterate communities reliably recognized and produced the same basic expressions observed in Western populations \citep{matsumoto_half-century_2022}.
Yet cultural differences still modulate how emotions are produced and represented mentally, where the core muscular patterns are conserved but expressive nuance and interpretation vary across groups \citep{jack_facial_2012, ruttkay_cultural_2009}.
Constructionist theories, on the other hand, argue that emotions are constructed experiences resulting from the brain's interpretation of internal and external stimuli \citep{barrett_solving_2006, lindquist_brain_2012}.
Behavioral findings support this second model as contextual cues routinely bias emotion perception, altering valence and categorical judgments even for prototypical expressions.
Neutral faces can be perceived as fearful or happy depending on preceding context, and using emotion labels can shape how facial cues are integrated before an explicit categorization occurs \citep{calbi_how_2017, brooks_conceptual_2018}.
These contextual effects are often automatic and difficult to suppress, which highlights the role of distributed, domain-general networks in emotion perception \citep{aviezer_automaticity_2011, lee_context_2012}.

Although emotions on real faces are expressed through subtle muscle contractions, emotions of virtual faces cannot be constructed in this way. 
One way to generate emotions on avatars is relying on the FACS-based principles that guide conveying emotions in virtual faces in a way that is analogous to emotions on real faces. 
Despite this, some emotions (e.g., disgust) remain challenging to render, and in some instances virtual faces elicit different recognition patterns than real faces \citep{garcia_design_2020, dyck_recognition_2008, de_paolis_perception_2015} which may influence the neural mechanisms associated with emotion recognition on virtual faces. 
On real faces, multivariate analyses sometimes recovers emotion-predictive patterns both locally and across distributed networks, but findings are often inconsistent \citep{kragel_decoding_2016, westgarth_systematic_2021, bendall_brief_2016}.
fMRI implicates a broad network, including medial PFC, amygdala, fusiform gyrus, superior temporal sulcus, and insula in emotion perception, with certain emotions preferentially engaging specific areas. 
fNIRS studies have likewise reported variable PFC responses to facial emotion, with some reporting increases, others decreases, or null effects in oxygenated hemoglobin (HbO) depending on task parameters and emotion type.
These findings highlight the need for more nuanced investigations of how different emotional expressions are represented in the brain.

Virtual faces can be just as effective as humans in conveying emotions, particularly when their facial muscle movements are clearly depicted \citep{hortensius_perception_2018}.
There is growing evidence that perceived authenticity and realism on faces modulate neural responses to emotions.
Since humans are highly attuned to perceiving real human faces, viewing avatars may engage different perceptual and neural processes, potentially leading to altered brain activity \citep{de_borst_is_2015}. 
For example, fearful human faces elicit stronger responses than fearful avatar faces in regions such as the superior temporal sulcus and insula, stronger theta and alpha EEG responses to avatars under certain conditions \citep{kegel_dynamic_2020}.
Using the same stimuli, it was found that avatar faces elicit significantly stronger reactions than real faces for theta and alpha oscillations \citep{sollfrank_effects_2021}.
This underscores the need for targeted comparisons of how a comprehensive range of emotions and kinds of faces are represented in the brain. 

fNIRS, EEG, and fMRI have previously been used to examine similar neural properties between emotions on real and virtual faces, however, a tool like fNIRS is well-suited to explore the same research question because it has reduced susceptibility to head movement artifacts and therefore is better for more naturalistic conditions, such as social interactions. 
fNIRS is also ideal for detecting activity in prefrontal regions which are strongly associated with emotion perception \citep{westgarth_systematic_2021, bendall_brief_2016}. 
Few studies have examined the interaction of facial emotion and avatar realism with fNIRS; existing studies with avatars focused on either a limited set of emotions \citep{kegel_dynamic_2020}, or none at all \citep{schindler_differential_2017}. 
To address this gap, we used fNIRS to compare localized cortical activity and functional connectivity elicited by static real and virtual faces across the full set of basic emotions, enabling a controlled comparison of face type and emotional content. 
We hypothesised that (1) virtual and real faces will produce distinct activation and connectivity profiles, and (2) different emotional expressions will evoke unique neural signatures. 
Given the exploratory nature of this study, we do not prespecify detailed directional predictions. 
Instead, this study aims to map how emotion category and face realism jointly shape cortical activation and network dynamics, informing the design of emotionally expressive avatars for education, clinical, and human-computer interaction contexts.